\documentclass[10pt]{article}

% --------- Minimal page and spacing setup ---------
\usepackage[margin=2cm]{geometry}
\usepackage{setspace}
\setstretch{1.05} % tighter than default
\usepackage{parskip}
\setlength{\parskip}{0.4em}
\setlength{\parindent}{0pt}

% --------- Math and Fonts ----------
\usepackage{amsmath,amssymb,amsfonts}
\usepackage{mathrsfs}  % for script letters
\usepackage{eucal}     % elegant calligraphic fonts
\usepackage{microtype} % better spacing

% --------- Section Style ----------
\usepackage{titlesec}
\titleformat{\section}{\large\bfseries}{\thesection.}{0.5em}{}
\titleformat{\subsection}{\normalsize\bfseries}{\thesubsection}{0.5em}{}

% --------- Header (optional) ----------
\usepackage{fancyhdr}
\pagestyle{fancy}
\fancyhead[L]{\small Notes}
\fancyhead[R]{\small \today}
\fancyfoot[C]{\thepage}

% --------- Document Starts ----------
\begin{document}

\section*{Categories}

A category is an abstraction composed by 
\begin{itemize}
	\item A class of objects denoted by $\text{Obj}(\mathrm{C})$. Objects on a \textit{category} $\mathrm{C}$ are the members of the \textit{class} $\text{Obj}(\mathrm{C})$ that is, the collection of all \textit{objects} of the \textit{category}. It is simply the \textbf{universe} of objects for the category $\mathrm{C}$. \textbf{Category} theory abstracts away from internal elements and focuses only on \textit{objects} and \textit{morphisms}.
	
	\item \textit{Morphism} are continuos maps and one must be careful on recalling that \textbf{homomorphisms} are \textbf{structure preserving} whereas in category theory there isn't any algebraic structure. They are just maps and every homomorphism is a morphism, but not every morphism is a homomorphism.
\end{itemize}	
\textbf{Example:}\\
\vspace{0.5cm}
In $\text{Obj}(\mathrm{C})=Set$

\end{document}
